On the Theory of Errors of Observation - Examples


BUILD NOTE (added as examples are made)

The page is now assembled from parts, as the Probability of Induction page is:
run ./build.sh in this folder to write index.html from src/. One part file per
example, so an example can be edited without touching the text. ./build.sh
--check says whether index.html is behind src/.

Examples are numbered in the order they are built, and the number appears in
the left margin beside the passage that opens it. Built so far: 1 and 2, both
in "The Theory of Observations". The examples for "Probability and Relative
Numbers" are not built — that section is still marked "##1 / Text /
Interactive Diagram" with no spec, so I have left the numbers free and started
at the first section that has one. If you would rather the numbering follow the
paper rather than the order of building, say so and I will renumber; it is a
one-line change per example.


Genearl:

always when possible have the examples coded into values or formulae in the text, clearly colour coded. 

Peirce uses esoteric notation here, so the standard should be to have a modern formula appear beneath his with the same colour code, each should then have = [version with the numbers filled in, colour coded] and then = [the answer]. if all goes well the answers should be the same... 

In general, we should let the text do the explaining, and let the graphic help make sense of the text. 

Avoid any sort of AI speak: its not x, its y, and that matters. etc etc. 

When appropriate to give additional context to the text or how the example relates., do not be rhetorical or pompous, just explain it simply. 


I have been giving descriptions to the examples. when you build them, please give them a number and place that number in the margin, like in the Probability of induction html. also include it in this document. 

for many of these, i want you to provide a graphical demonstration. i will then suggest any tweaks. For graphical demonstrations, we want the user imnput to happen as much as possible through sliders that change a graphic and corresponding numbers in real time. 

as much as possible, you should try to tie new examples in to ones that the user will be familiar with, rather than trying to show it in a new way, show how this example is really like a previous one, but with this or that aspect of it added, removed, changed, etc. 


WHen you are done, do not move the eamples or anything off of this list. Add to the bottom of each example a sort of progress log that i can read, so i can update as we go. 


# Probability and Relative Numbers

Should have graphics for all the main bits, the universe under consideration,, how different formulae pick out different bits of it, how the dividing works...  ideally it would be nice to also show the modern notation beside it. 
it ends with Bayes theorem. 

will have to make a note here about Gauss and the different ways of justifying the ToE. Peirce's set up looks Laplacean / bayesian, but the rest of the paper shows he is in fact gaussian / frequentist about it. 



##1

Text





Interactive Diagram



# The Theory of Observations

## Gaussian Theory of Errors 

"We now pass to the theory of observations. An observation gives us the value of a certain quantity which is connected with an unknown quantity in such a way as to be partly dependent on the latter value, and partly on accidental circumstances, not capable of being separately taken account of."

can say that this is the Gaussian Theory of Errors  and now we put this in terms like this: Oi = T + \epsilon_i 

## 1 — The quantity observed and the quantity wanted
Status: awaiting approval

### Text
"We now pass to the theory of observations. An observation gives us the value
of a certain quantity which is connected with an unknown quantity in such a way
as to be partly dependent on the latter value, and partly on accidental
circumstances, not capable of being separately taken account of."

### Suggestions
- Is the transit-instrument framing right this early, or should the first
  example be neutral about what is being measured?
- [from the build] The three-row statement — words, then Peirce's symbols, then
  the numbers — is the pattern I propose to reuse for every formula in the
  paper. Say if you want it the other way up, or the modern line added as a
  fourth row rather than as a footnote.
- TO EDIT: the opening paragraph of the example is a placeholder of mine and is
  yours to rewrite. It is the only prose left in the example, and it carries the
  colour key for the whole paper — unknown quantity green, accidental
  circumstances red, quantity observed blue — so whatever replaces it should
  still name those three things; the three coloured spans are in the markup and
  want keeping. The transit instrument as the running example is open too. It is
  marked with a TO EDIT comment at the top of src/06_ex1.js.

### Awaiting approval
#### the statement at the top, the cumulative cloud, the residual
Three colour-coded phrases in Peirce's own sentence light when the example
opens. Sliders set T and the reach of the accidental circumstances; buttons
observe once or fifty times. A number line shows T as a green rule and each
observation as a dot, the newest bracketed back to T and labelled with its own
epsilon; underneath, the series piles into a histogram. The statement sits at
the top, above the controls:

    the quantity observed = the unknown quantity + the accidental circumstances
             O_i                     T                        e_i
            4.911                  5.000                    -0.089

The = and + are set once, in the top row only. The cloud is cumulative: every
observation is kept and drawn faint at a vertical position fixed when it was
made, so the pile darkens where the observations fall thickest — measured on the
built page, the dense middle reads about 30% darker than the sparse edges.

A checkbox hides the unknown quantity; nothing in the data changes when it does,
which is the observer's actual position. A second checkbox shows the mean as a
gold dashed rule, and with the unknown quantity hidden the third term becomes
the residual — the bracket is measured from the mean, labelled r, and the
statement reads O_i = x-bar + r_i in gold rather than green and red. With both
shown, both brackets draw, and the gap between them is the difference between an
error and a residual.

All closing prose cut, at your instruction. The example ends at the plots.

### Completed
None yet.

## 2 — The law of the facility of errors
Status: awaiting approval

### Text
"These accidental variations are, however, in all cases subject to a statistical
law, so that (observations of a certain kind forming the limited universe, X
being the unknown quantity, Xi the quantity observed) the quantity,"
... "The special form of the function phi is called the law of the facility of
the errors."

### Suggestions
- The strip is fixed at 0.10 wide. Should it be a slider, so that d-epsilon
  getting small is something you can watch? This matters more now that the strip
  is carrying Peirce's d-xi.
- "Two observers run together" anticipates the splitting-the-universe section a
  long way ahead. Keep it here, or hold it back?
- [from the build] The limited universe is only implied — it is whatever kind of
  observation is selected, and changing the kind clears the shots. Worth making
  explicit in the graphic, or does the radio button carry it?

### Awaiting approval
#### phi as a relative number, folded in from the passage above
Four kinds of observation, all named by Peirce: transit observations (normal), a
star out from behind the moon (skewed right, no mass to the left), a coarse
instrument (uniform), and two observers run together (bimodal). The marksman is
on the left and the law on the right — the same fact drawn twice, since the
horizontal scatter of the shots is what the curve describes. A strip of width
d-epsilon runs across both.

The passage between examples 1 and 2 is where phi comes from, so it now opens
this example too. Peirce defines it as a relative number, and the example was
already computing that ratio without naming it, so the statement is two-sided —
a tally on the left, a law on the right:

    the shots in the strip, out of      =     the height of the law there,
    every shot fired                          times the width of the strip

    [xi_Xi , x_x] / [x_x] . d-xi              phi(e, x) . de

    24 / 400 = 0.0600                         0.807 x 0.10 = 0.0807

The bracket notation uses the same mirrored inverted comma as the article text.
The readout gives both numbers and says they close on each other as more shots
are fired, which is the whole content of the claim. A line underneath notes that
the left side is a tally and nothing else, that the limited universe here is
every shot this marksman fired, and that since e = xi - x the two differentials
are the same width.

Example 2 now opens from three passages; only the first takes the margin
number, which needed a change to the shared scaffolding.

Two layout bugs fixed on the way: a note spanning every column of the statement
grid was dumping its whole width into the operator column (337px), and the
three-column variant outranked the narrow-screen rule so it never collapsed on
mobile.

### Completed
None yet.

## Need for a law of error

"Except so far as this law is known, an observation can afford us no information whatever."

have the same marksman / gaussian analogy, but this time take away the target / true value. 

show how the same set of data can be explained by any number of combinations of other distributions or multiple shooters. 

might need something about how the empirical distribution is not a law of error. 

this shows how we need to have warrant for believing that we are working in the same universe (one shooter / observer)  and that we know the law governing those shots / observtions. 

please suggest ideas for how to visualize the stronger point that w/o a LoE there is no information whatever. 

## Conditions

We should visualize the following (and what it would mean to break them):

It has but one value for each set of values of its variables.
Its value is always positive and less than unity.
It vanishes when ε is infinite.
Its integral relatively to ε from −∞ to +∞ is always unity.

this should be graphical, but also show the function itself in action, so that you can see the values change. see what happens to the whole if one or another is not satisfied. 

we shoudl say what the modern name for these assumptions is. is this zero conditional mean? we should say. 

## x-z

Show this:

"The value of this probability is for any particular kind of observations an arithmetical function of ε and ξ, which we may write ψ (ε, ξ) dε."

then this:

"The probability that the unknown quantity X has the value x without reference to observation is [xx]dx. This is in any case a function of x, which may be written Ψx . dx."

and this:

"The probability that the quantity given by observation Ξ has the value ξ, without reference to the value of the unknown quantity, is [ξΞ] dξ. This an arithmetical function of ξ, which may be Φξ . dξ."

and this:

"If Φξ, Ψx, and φ (ε, x) are given, then we can obtain ψ (ε, ξ) thus:"

(this final one should clearly show how it is the result of the others, letting the user see it expanded or contracted, with colour coded terms, a modern equivalent, and a graphical example)

## Series of observations. 

"Suppose a number of independent observations to be made. Then we shall have a series of functions—"

this should let the user set how many series, which would add additional functions to the list we can give them numbers, and these shoudl carry into the next bit too. 

## x

"The probability Ψx . dx antecedent to all observations will be simply dx, and, therefore, the factor Ψx may be omitted in the above expression."

not sure how to show this. 


## Observation teaches a function (least squares)

This paragraph "It will be perceived that observation never gives us to know a number expressing the value of the unknown quantity, but only a function expressing the probability of each value."

This can use the above machinery to show that we never get T=O but always Oi = T +ei, and the ei bit is sort of where the function comes in (as it is a function of the type of observation).

at first let them select from  / create a whole host of distributions

have a button that sets it to the very comprehensive case, and show how it lets you describe the whole thing in terms of mean and PE. colour coded etc. 


##MR. Crofton

give us two observations of X. so X1 = zeta 1 = number and X2 = zeta 2 = number, just numbers to carry through. (user could change them on slider. ) let the user add a third - 5th one too. interacting with this bit :"If in place of two we have n observations, the probability that the sum of all will"

then carry these through the math, colour coding and showing any appropriate graphics. 

we will want to demonstrate this bit "is half the probable value of the square of this error, and so on. "


##M is large

"he probable value of the error is often less than the probable value of its square, owing to positive and negative errors balancing. But the co-efficients which involve the cube and higher powers of the error may often become insignificant." to "Consequently if M is sufficiently large, the higher co-efficients may be neglected."

we will want to make the formulas live and show peirce's point graphically




## Galton Board

"The manner in which sensation and volition are propagated through the nerves is unknown, but it must be by some very complicated process, because it is very slow compared with the rate of propagation of sound. It is, therefore, probable that there may be variations of the rate of passage, and that the velocity through each small portion of the whole length of the nerve is to some extent independent of the velocity through the other parts. If this is so, the whole time of propagation would be subject to a variation, the probable values of whose cube and higher powers would be insignificant. However this may be, it appears to be a fact that in all carefully-made observations, the error of which is due to the inevitable inaccuracy of the action of the human nerves, the probable values of the cube, &c., of the errors are very small."

Have a graphic that is like a horizontal galton board. have a slider for how complex the process is. at its simplest, it is just a stimulus and like a single line to the brain. as it gets more complicated, it is more of a tangled branching set of lines, and as the stimulus hits, you can watch it travel through the lines in a sort of complicated way (each one of the lines carries it at a different speed. 

connect this to the cube etc stuff, with corresponding graphic. 


##CLT

"We thus reach the fundamental formula of the method of least squares."

we need to do two things here. first, show how it relates to modern CLT, which should be next to it.  show how h relates to sigma. have sliders that really show how "h and E are quantities which depend upon x."


##FUnctino not a number again

"It is not the object of this paper to explain that method itself. Only it may be remarked that in the deduction of it which is usually given, it is assumed that what we wish to obtain from observations in any case (whether the method of least squares is applicable or not) is the most probable value of the unknown quantity. This is not the case, for there is but an infinitesimal probability in favor of any one value. Suppose that the cause of error in observing the place of a star were a nearly simple oscillation of the image about its mean point. Then the most probable errors would be the extreme ones. But we should much prefer to assume as the value the probable or mean value than the most probable value, which would be removed the furthest possible from that value. The conception of the matter is this. What observation has to teach us is a function, ψ (ε, ξ), not a mere number. But in cases to which the method of least squares is applicable, this function is completely determined by two numbers, which are the value of the unknown quantity derived from the application of that method and the value of its probable error.2"

we will sort of have this from above, so we can build it out to make these points. 

##Weighting

"It is assumed in treatises on least squares that Φξ . dξ, or the probability (without regard to the value of the unknown quantity) that the quantity observed will have a value between ξ and ξ + dξ, is equal to dξ. When this is not the case it is only necessary to weight each observation by dividing by Φξ."

we will want a graphical demonstration of this. 

show this:
"If the probability of error does not follow that law which the method of least squares supposes, that is to say, if the probability of

the error x in the mean of a large number of observations is not equal to 
h
ö
G−h²x², where h is a constant independent of x, then it must at least be of this form if h be considered to be a function of x. "



##weighting again

"Now, h² is the weight which has to be assigned to an observation in the application of the method of least squares; and therefore when this method does not apply in its unmodified form, it is only necessary to find what function of x, h must be in order to give the required law of facility of errors, and then proceed according to the method of least squares, after having weighted each observation by h². "

we will want to show this using the example he gives just after. i.e.,:

""Let the equation which represents the facility of error be plotted, the error being taken as abscissæ, and the probability of that error as ordinate;

then plot on the same diagram the curve y = 
1
ö
 G−x². Let us suppose, then, that a certain value of x, y, is A times as great for the actual curve of errors as it is for the normal least-squares curve. Now, if h be increased in the ratio A, the ordinates will be increased in this same ratio, and the abscissæ will be diminished in the inverse ratio so that the area of the curve is preserved. But if the ordinates at any one point are to be increased in the ratio A, then the abscissae at that point must be contracted in the inverse ratio, so as to preserve the area; so that if we had a function fx such that Dx fx = 
1
a
, then the probability of this function of the error would follow the law y = 
1
ö
 G−(fx)², and consequently 
1
(Dx fx)²
 is the weight which has to be assigned to any such observation in applying the method of least squares. "
 
 
 this graph should appear after that paragraph. make it visually clear how the weight depends on the value of the error.  and show this approximate step and its importance:
 
 "it will be necessary first to make an approximate solution of the problem in order to get an approximate value of the error from which to determine the weight, so that when the method of least squares is not applicable in its unmodified form, an approximate method must necessarily be resorted to."
 
 ##Ascertaining the law of the facility of error

"In any case, if the error is compounded of an infinite number of infinitely small errors, or approximates to being so, then the law of facility of errors is of that general form which the method of least squares prescribes, and nothing remains indeterminate excepting the value of h." 

hve a graphic showing this. could be a sort of galton board or something. 

## Table

"I copy, for example, from Chauvenet’s Astronomy the following table, taken from Bessel’s Fundamenta Astronomiæ, showing the errors made by Bradley in observing Sirius and Procyon:"
 
 lets have a graph of this table.
 
 ##Fechner
 
 "Fechner, in his Elemente der Psychophysik, has, in connection with his psychophysical law, discussed the applicability of the method to cases of sensation generally such as the estimate of the relative weights of two masses, and he finds that if v be the energy of the force which produces the excitation of any nerve, then if log v be considered as the observed quantity, the errors of the observed quantity will follow the law of least squares. Strictly speaking, the law of least squares recognizes the possibility of any error positive or negative, however great, although the probability of indefinitely great errors will be indefinitely small. "
 
 Use the nerve galton board again, but showing how the results at the far end will give a normal distribution. 
 
 Link to Fechner's Lab
 
## Occultation

"It occurred to me that in the case of the emersion of a star from an occultation, since it was impossible to strike the chronograph-key too early, while it might be struck indefinitely too late, the law of least squares could hardly apply, and I have made some experiments upon this point, which I will narrate in detail at the end of this paper, merely remarking in this place that the divergence from the theoretical law proved to be insignificant. "

Have two views of a transit observation, one where the star is seen coming and one where it appears suddenly from behind a moon. have a button to run one simulation, which releases on star on each and has an imaginary person "observe" it, making a line where their observation time is. have this line also appear as a data point on a chart below the observation screens. Have another button that observes like 100 or 1000 transits, which would fill in the regular one as normal and the occultation one as having only a skew to the right. making it clear that his is peirce's supposition. 

there could also be a mode that lets the user try it out but the stars would have to move pretty fast for it to be any fun. it would clear the other data and let them try for themselves with the same mechanism. 


## Large errors in chemistry. 

Show the point made in the paragraph up to "The method of considering such errors is treated by Dr. Bremiker, in the preface to his Tabula Logarith- morum Sex Decimalium, a work to which the attention of chemists ought to be drawn."

let the user add a large error in and all that. 

if you can find or understand the theory peirce refers to, lets display it as well. 

## Not a priori

Show this:

"When the law of facility of errors cannot be deduced a priori from the consideration of the causes to which it is due, a large number of experimental observations must be made upon a known quantity in order to find in what manner the errors vary, or the same series of observations may be used both to deter- mine what the value of the unknown quantity is and also what the law of the variation of errors is. Thus, in the method of least squares, h is to be determined empirically, and the common way of doing it is to use the actual observations by which the unknown quantity is determined to determine also the value of h. In doing this, it should be remembered that a precise and trustworthy determination can only be obtained from a large number of observations. "

## New unkonwn

show this "his proce- dure amounts, in fact, to adding an additional unknown quantity, a very obvious fact, and yet one which is habitually overlooked. Encke, in his Astronomisches Jahrbuch for 1834, has given the most com- plete formulas that I have anywhere found for determining the value of h, as well as its uncertainty. "

##Combination

"These formulae require correction on account of the circumstance just mentioned. When it is necessary to combine, by least squares, observations of different orders of precision, they are weighted proportionately to h2. If we have two series of observations, one of which is as accurate

as you please, and the other as inaccurate as you please, a better result than that which the most accurate series of measures gives can always be got by combining with it the least accurate series, provided the proper weights are given to the two series. ... The least accurate series of measures offers certain facts, which may be used as premises, and it cannot be that if those facts are properly used they leave us in greater ignorance than we were before. Additional facts must increase our knowledge, if the proper inferences are made from them, and, therefore, an additional series of observations, if it have any weight at all, must, if its proper weight be assigned to it, improve the value of the unknown quantity. "

have two distributions, one for each series, and then a combined one. let the user change the weights and see how the combined distributino changes. We can note that CSP's method for showing this is coming up. 


## Combination 2

"On the other hand, when it is considered that there is an uncertainty in the value of h, it may be that if the two series of observations differ greatly in accuracy, and the value of h is not determined with much precision, it may be better at once to take the result of the best series of observations than to combine the two series with the best weights that we are able to give."

show this in a similar way, letting the determined value of h cahage too. the user sohuld be reminded visually of the earlier example where the value of h is determined empirically. 

## formula

"xi be the value from any set of observations;
εi the mean error of this set;
giεi the mean error of εi;
wi the true weight;
E the mean error of weight one; and
x the truly weighted mean.
x = 
Σi(wixi)
Σiwi"

usual color coding etc. 

tie these together using graphics the user is now familiar with. 


we will want a graphic for each stage of the math until "Some writers upon the subject have wished to assign a smaller weight to those observations which differ largely from the mean than to those which come close to it."

we should have two series of observations. the user can set some parameters of them. then we have peirce's math's guess as to whether it would be better to combine. a button reveals the true value and we see if CSP was right.  (If that is what he is doing... to be honest i am not sure. so you let me know what he is up to here. take this part seriously as it is the core of the paper and something hard for me to check)


## H is a statistical quantity. 

show how h cannot belong to a single observatoin, but only and lawys to a series. 



## Splitting the universe

"They have reasoned as if h, or the precision of an observation, were something which belonged to a single observation; whereas, in fact, it is a statistical quantity alto- gether, and belongs only to an observation as a member of a certain series. We may have a large series of observations for which h has a certain value, and those observations may perhaps be separated into two series on some principle or other for which h shall have two different values, and if this can be done there is an advantage in doing it. It is, in fact, limiting our universe. In probabilities generally h is a mean or probable quantity for a series of observations, and if we can divide our universe into two parts, getting different values of h, it will be an increase of knowledge to do so. For example, suppose that some of the observations were taken under one set of circum- stances and the rest under another set of circumstances. That would afford a principle upon which the observations could be distin- guished, and if the value of h for the two sets turned out different, it would be an advantage to separate them and to give them different weights. "

show how a single set can be splipt up, and how the different distributions can be recombined to advantage. this shoudl recall grahpics for the universes from the probability section. 



## Peirce's Criterion

Let the user apply Peirce's Criterion

"f observations are to be rejected, exact criteria are necessary to determine upon principles of probabilities in what cases they should be rejected. The criterion of Professor Peirce is the only one which conforms rigidly to those principles, and, indeed, I am not aware that it has been attacked upon the ground of not conforming to the principles of probabilities, although it has been attacked on the ground that no such criterion should be used, and that no observation should be rejected except upon principles of guesswork, for that is what it amounts to to say that none but obvious mistakes should be rejected."

It sohudl show the data graphically of course, show how things change for each observation being considered the error, and make red those considered to be outliers. 

Details here: 

Gould, B. A. 1855. “On Peirce’s Criterion for the Rejection of Doubtful Observations, with Tables for Facilitating Its Application.” The Astronomical Journal 4 (April): 81. https://doi.org/10.1086/100480.
Peirce, Benjamin. 1877. “On Peirce’s Criterion.” Proceedings of the American Academy of Arts and Sciences 13: 348. https://doi.org/10.2307/25138498.
Ross, Stephen M. 2003. “Peirce’s Criterion for the Elimination of Suspect Experimental Data.” Journal of Engineering Technology.


 "Experience has shown that the errors which this criterion rejects are almost precisely those which a person of sound judgment would pronounce to be obvious mistakes, but some other criteria have been proposed, which are confessedly inex- act, but which have the advantage of involving less calculation, but these are no better than the unaided judgment of an experienced person, and in some cases not so good."
 
 let the user guess at the outliers, and see if they are right. 
 
 do you know of any of the other methods? they could be incorporated if so. 


#Experiments

## Accuracy

Show this with a simple simulaiton. even with precision capped at 1ms, the mean can get better. 

"e instrument ought, theoretically, in the mean of a large number of observations, to give a much closer result than 0.001 second; for when the first event occurs, and the arm is thrown into the moving crown-wheel, it probably will not strike exactly in any catch, but will strike on the inclined side of a tooth. If it strikes on the forward side of the tooth,

the hands will be carried forward by the fraction of a thousandth of a second more as the arm glides down this side to the bottom of the catch. But if it strikes on the back side of the tooth, the hands will be carried relatively back the fraction of the thousandth part of a second as the arm glides back to the bottom of the catch. Now, if the top of the tooth is midway between the bottom of two catches, it is equally likely to be carried forward or back. The same thing occurs when the second event happens and the arm strikes upon the fixed crown-wheel. An error in the marked interval will thus result, which error may amount to −0.001 second in the extreme, or to +0.001 second, and any one error between these limits is as likely as any other; consequently, these errors, being entirely incidental and inde- pendent of one another, they will balance one another in the mean of a large number of observations, and thus a much higher degree of accuracy may be reached. "

##Experiment

Copy the occultation simulation from above. let the user try it. have. button to simulate 500 and give a graph like those Peirce found. 

## 3 — The twenty-four days: the tables and the plate
Status: awaiting approval

### Text
"Five hundred observations were made on every week-day during a month,
twenty-four days' observations in all. The results are given in the accompanying
table, and are also shown upon plate No. 27."

### Suggestions
- The pencil ("Draw it myself") draws left to right only and does not yet let
  you go back over a stretch. Enough, or should it be a proper freehand tool?
- The plate is drawn at one scale per panel, each centred on its own day. Peirce
  drew them on a common horizontal scale, which is what makes the drift across
  the month visible down the column. Say which you want. Easier to judge now
  that his own plate is a button away.
- Nothing yet uses the tables to make the paper's own point — that the range of
  errors narrowed after the twelfth day. That wants its own example, on the
  sentence that says so.
- [from the build] Should the discrepancy note live in the example, or only in
  these notes? At the moment the tables tab says when a day does not sum to 500.

### Awaiting approval
#### REVISION 3 (your notes): the plate's two buttons, the legend cut
- The plate carries both buttons at all times — "The redrawing" and "Peirce's
  own plate" — with the current one marked, instead of one button that swapped
  its own label. It opens on the redrawing. The mean-curve and reset controls,
  and the two sliders, still appear only on the redrawing, there being nothing
  for them to drive on his.
- The line under the calendar is cut. The question I had raised with it — why
  the 2nd to 4th and 11th to 13th of July are missing — goes with it.

#### REVISION 2 (your notes): the plate's scales, the day's table, the pencil out
- Every panel of the plate now carries its own scale: the four verticals are
  labelled with the millisecond they stand at, set under the baseline. The
  canvas is taller to make room. They are in milliseconds, to agree with the
  tables and the one-day tab; Peirce labelled his in seconds (0.15, 0.20 ...)
  and only on the top panel of each column, since his panels share a scale.
  Say if you would rather have it his way.
- On the one-day tab the table has moved up beside the buttons — it is now in
  the same column as the plot, directly under it, so the calendar, the sliders
  and the buttons sit alongside both. It scrolls within 300px.
- "Draw it myself" is gone, and the pencil machinery with it.
- The Plate 27 figure that sat in the article above this block is removed. The
  plate is a button away inside the block now, so it was there twice.
- The mean curve has a slider of its own, from 4 (bouncy) to 48 (very smooth),
  default 22. It sits beside the smoothing slider on both the plate and the
  one-day tab, so on the plate it redraws all twenty-four mean curves at once.
  "Reset to Peirce's" restores both sliders.

#### REVISION 1 (your notes): at the foot of the page, and a calendar
- The whole thing now sits open at the foot of the page and is not behind a
  trigger at all. It is no longer an example container: it is a plain block set
  off by a rule, with the numeral 3 still in the margin, built at load. The
  "Five hundred observations" sentence is plain text again.
- The day is chosen from a calendar of July 1872 with August 1-3 beside it,
  instead of a dropdown. July 1872 began on a Monday. Days with observations are
  clickable and carry their roman numeral and total in the tooltip; days without
  are dimmed, which puts the shape of the month on view — every Sunday absent,
  and the 2nd to 4th and 11th to 13th of July.
- The plate has a button that swaps between the redrawing and Peirce's own
  engraved Plate 27. On the redrawing the smoothing slider sits above it, so
  moving it redraws all twenty-four panels at once and the month tightens
  together; the mean curves toggle with it. The slider is hidden on his plate,
  there being nothing to drive.
- "Reset to Peirce's" on both the plate and the one-day tab: eight passes, mean
  curve on, everything else off.
- The chosen day carries between the tabs.

#### three tabs: the tables, the plate, one day
Data from Koenker's transcription (Koenker's data/MiM/data), parsed exactly as
his ReadPeirceData does, and held in the shape Peirce printed it: nine
column-pairs per day, each read downwards, with a value smaller than the one
above it treated as an abbreviation for the hundreds carried down. 13 KB
embedded in the page, so nothing is fetched at run time.

THE TABLES. Peirce's own layout, reproduced to the digit. His abbreviation rule
turned out to be readable off the printed page: the time is set in full at the
head and the foot of each column, wherever the run of milliseconds breaks, and
at every multiple of ten; elsewhere the last digit alone. Checked cell by cell
against pp. 138-139 — the first row reads 158 1, 348 0, 389 1, 430 0, 471 1,
512 1, 553 0, 593 1, 633 2, and the first column reads 158, 240, 261, 277, 296,
312, 3, 4, 5, 6, 7, 8, 9, 320, 1, 2, 3, 4, 5, 6, 7, both exactly as printed.

THE PLATE. All twenty-four panels, two columns of twelve, roman numeral and date
at the left of each, baseline and four verticals as on Plate 27. Each panel
carries the figures smoothed eight times over and a mean curve.

ONE DAY. The smoothing is put in your hands. Peirce says the curve was "smoothed
off by the addition of adjacent numbers in the table eight times over", which is
convolution with [1,1] eight times, so the slider runs 0 to 16 passes and is
labelled at 8 as his. At 0 the raw figures show as spikes. The sums are divided
back by 2^passes so the vertical scale stays comparable. Then the four overlays:
a mean curve by eye, a pencil so you can draw your own — his smoother curve is
"a mean curve for every day drawn by eye", not a fit of any kind — a Gaussian of
the same mean and spread, and clear.

#### discrepancies found in the data, as asked
1. Peirce says five hundred observations were made on every week-day. The
   printed tables sum to 500 on only two of the twenty-four days (the 14th and
   the 19th). The whole month comes to 11,803, not 12,000 — 197 short.

2. The twelfth day, July 20, sums to 396: 104 short, where no other day is out
   by more than 11. This is not a transcription slip. I summed the printed table
   on p. 149 by hand and it comes to 396 as well, and that day's table is
   visibly shorter than the others — 18 rows where the eleventh day has 20, and
   it shares a page with the thirteenth. Whatever happened, happened before the
   tables were set.

3. Two days run over: the seventeenth by 7 and the twenty-first by 2.

4. Koenker's Day01.txt has two stray periods, "1." on line 9 and ".0" on line
   15. R's scan() reads them as 1 and 0, so his analysis is unaffected, and the
   parser here does the same; they are transcription noise, not data.

### Completed
#### the mean curve's default width
22 is right. [22 seems good.] The slider stays, defaulting there, and "Reset to
Peirce's" returns to it.

#Graphs and table

have the graphs and table both ready, but part of a set of tabs. it can display tables, or graphs, or the tables and graph for a particular day. 

##Table

Recreate Peirce's tables of data from Koneker's data*.  Copy the style from the text. 

(while doing so, check for any issues or discrepencies)

##Graphs

reproduce Peirce's graphs (including their style and layout) using Koenker's data*.
 
 Let the user set the level of smoothing and drawing curves using the same sort of app as in the three methods html in the folder. 
 


* Koenker's data is available in the folder of the same name and the smoothing procedure is described in three-methods html. 